\section{Spanning by the Poincar\'e map}\label{sec:spanning}

\begin{theorem}[Spanning theorem B]\label{thm:spanning}
The dimension \(h\) of the target space \(V\) is positive. There is a finite
set \(T\subset\bar G\) such that
\[
  |T|\dim S\leq 5h\log(1+2b^2),\qquad
  \dim\sum_{g\in T}\mathfrak p_g S\geq (1-4/b^2)h.
\]
\end{theorem}

To control the span of the Poincar\'e series, we estimate inner products between
translates by the finite level group. We first work with the full source space
\(S_0\) and then impose the ramification condition.
Let \(\mathcal H\) be the space of forms of principal congruence level on \(Y_0\),
with the unitary action \(\rho\colon\bar G\to U(\mathcal H)\). The orthogonal
projection \(e_0\) onto \(V=e_0\mathcal H\) commutes with this action.
Put \(D_X=\dim S_0>0\). We equip \(S_0\) with the inner product coming from
the double cosets and Bessel coordinates of the formalization. The norm
\(\|\cdot\|_{\mathrm{HS}}\) is the Hilbert--Schmidt norm.

\subsection{Unfolding and local bounds for the second moment}

Let \(\deg_B\) be the product of the formal degrees of the Bessel systems at
the infinite places, and let \(\mathrm{mass}_r\) be the mass of the genus of
the standard Hermitian lattice of rank \(r\). The mass is the sum over
isomorphism classes of lattices, weighted by the inverse orders of their
automorphism groups. For the target space we have
\begin{equation}\label{eq:span-dimension-mass}
 h\leq |\bar G|\deg_B\,\mathrm{mass}_r.
\end{equation}
Equality of multiplicities identifies \(h\) with the dimension of the
corresponding projected space on the definite space \(W_m\).
\leanref{targetDim_eq_finrank_goodForms_spaceW}
On \(W_m\) the classes at principal congruence level have trivial
automorphism groups. Their number is therefore the mass \(\mathrm{mass}_r\)
at maximal compact level times the level index \(|\bar G|\). By Weyl's dimension
formula, the product of the dimensions of the algebraic representations
attached to each class is \(\deg_B\). Hence the dimension before projection
is the right-hand side of \eqref{eq:span-dimension-mass}.
\leanref{finrank_principalForms_spaceW}

Let \(J\) be the principal congruence level and \(\nu\) the covolume of the
real unipotent lattice, which is independent of the class. In the coordinates
of the formalization, scale the Poincar\'e map by \((\operatorname{vol}(J)\nu)^{-1/2}\)
to obtain \(P\colon S_0\to\mathcal H\), and put
\[
 P_s=P^*\rho(s)P \qquad(s\in\bar G).
\]
Let \(c,c'\) be classes of the source space on \(U(X_0)\), let \(h_c,h_{c'}\)
be the images of their representatives in \(U(Y_0)\), and write \(W\) for the
Bessel function and \(\beta_c,\beta_{c'}\) for the class cutoffs. Unfolding the
inner product of the Poincar\'e series gives the \((c,c')\) block of \(P_s\)
as the entrywise complex conjugate of
\begin{equation}\label{eq:span-unfolded}
 \frac1\nu
 \sum_{\substack{\gamma\in U(Y_0)(F)\\
       h_c^{-1}\gamma h_{c'}\in K,\!
       \operatorname{red}(h_c^{-1}\gamma h_{c'})=s}}
 \int
    \beta_c(y)\beta_{c'}(\gamma^{-1}y)\,
    W(y)W(\gamma^{-1}y)^*\,dy .
\end{equation}
Here \(F\) is the maximal totally real subfield of \(E\), \(K\) is the lattice
stabilizer in the finite adeles, and the integral is over the product at the
infinite places. \leanref{hasAmplitudeBlocks_scaledPoincare}

Let \(R_0\) be the rational Bessel subgroup, and let \(\bar R_0\subset\bar G\)
be the reduction of its integral points. Restricting
\eqref{eq:span-unfolded} to \(\gamma\in R_0\) gives the unitary representation
\(r_\psi\) of \(\bar R_0\), extended by zero to \(\bar G\). Set
\[
  E_s=P_s-r_\psi(s),\qquad
  D_{\mathrm{id}}=|\bar R_0|D_X,\qquad
  D_{\mathrm{main}}=\frac{D_X^2}{\deg_B\,\mathrm{mass}_r}.
\]
Write \(\bar G=G_0\times G_1\), where \(G_i\) is the product over the chosen
places above \(p_i\), viewed also as a subgroup of \(\bar G\).
The calculation of the terms on \(R_0\) gives
\begin{equation}\label{eq:span-equivariance}
 P_t=P_1r_\psi(t),\qquad E_{ts}=r_\psi(t)E_s
 \quad(t\in\bar R_0,\ s\in\bar G).
\end{equation}
Indeed, the action of \(r_\psi(t)\) on the source reindexes the terms of
the unfolded expression. \leanref{besselTranslate_mul_left}

To bound the second moment on larger sets, group the terms in
\eqref{eq:span-unfolded} by their correspondences attached to the isotropic
part \(I\) of the Bessel frame. The correspondence of a term \(\gamma\) is
the subspace consisting of the pairs of images of \(y\) and \(\gamma^{-1}y\)
in \(I^\perp/I\), for \(y\) orthogonal to both \(I\) and \(\gamma I\).
For fixed classes, reduction \(s\in\Omega\), and correspondence, sum the
amplitudes of the truncated Bessel function. Let \(\mathcal L(\Omega)\) be
the sum of the squared norms of these sums divided by \(\nu^2\).
The terms included are those for which
\(\mathfrak d(\gamma)=\Delta\overline{\Delta}-N^2\Delta'\overline{\Delta'}\ne0\),
where \(\Delta\) is the determinant of the frame pairing matrix and
\(\Delta'\) is its preceding leading principal minor.
This is the normalized sum over links.
\leanref{linkSecondMoment}

We use the primitive \(5040\)-th root of unity in \(E\) and the Weil splitting
character chosen in the formalization. The Bessel system and level constructed
from admissible data satisfy the size and level hypotheses of the following
three estimates. \leanref{filling}

\begin{enumerate}
\item
If \(\Omega\) is disjoint from the cosets of \(\bar R_0\) obtained by
multiplying by nontrivial roots of unity, then
\begin{equation}\label{eq:span-cutoff}
 \Bigl(\sum_{s\in\Omega}\|E_s\|_{\mathrm{HS}}^2\Bigr)^{1/2}
 \leq \mathcal L(\Omega)^{1/2}+b^{-15}D_{\mathrm{id}}^{1/2}.
\end{equation}
Section \ref{sec:span-pair-count} bounds the cross terms in the sum of squares
and the error from truncation.
\leanref{sqrt_secondMoment_compl_besselS_le}\leanref{secondMoment_remainingTerms_le}
\item
Among the chosen places and their conjugates, take those at which the
span of the reduced frame has a nonzero vector in common with its image
under \(s\). Let \(\operatorname{Meet}(s)\) be the product of their norms.
If a positive integer \(z\) is at most \(\operatorname{Meet}(s)\) for every
\(s\in\Omega\), then
\begin{equation}\label{eq:span-intersect}
 \mathcal L(\Omega)\leq
       \left(\frac{b^{2f}}z\right)^k D_{\mathrm{main}}.
\end{equation}
Section \ref{sec:span-pair-count} counts pairs in the same correspondence
using the lattice in the intersection. \leanref{linkSecondMoment_le_of_le_meetNorm}
\item
Let \(A\) be a set of chosen places above \(p_0\). For each \(w\in A\),
specify a nonzero vector in the span of the frame at \(w\) or its conjugate.
Write \(\Gamma_A\) for the subgroup fixing these vectors.
Call a reduction \(s\) open if at some chosen place above \(p_0\) all
proper leading principal minors of positive order of the frame pairing
matrix are nonzero. Let \(\Omega_A\subset\Gamma_A\) be the open reductions.
Then
\begin{equation}\label{eq:span-open-input}
 \mathcal L(\Omega_A)\leq (1+b^{-4})b^{-6k|A|}D_{\mathrm{main}}.
\end{equation}
Section \ref{sec:span-open-proof} applies Siegel's estimate and Poisson
summation, deriving a factor \(b^{-6k}\) per fixed-vector condition from
the local mass. \leanref{linkSecondMoment_isOpenElement_le}
\end{enumerate}

\subsection{Counting with the common lattice}\label{sec:span-pair-count}

Put \(n=\varrho+1\), \(n_X=\#\{c\}\), \(\tau=b^{kf/4}\), and
\(n_p=|\bar R_0|n_X/N^{kf}\). Group the terms of the truncated Bessel function
into double orbits under the integral unipotent groups. A term
\(\gamma\) is near if there is a point \(y\) in the product at the infinite
places such that both \(y\) and \(\gamma^{-1}y\) lie in the region of reduced
points with bounded Bessel lengths. In each orbit containing a near term,
choose one of them as representative. A reduction \(s\) matches the left and
right unipotent characters if their values agree on integral unipotent
elements satisfying \(us=su'\) after reduction. If this condition fails,
the change of variables by such a pair multiplies the amplitude by a
nontrivial character ratio and forces it to vanish.

The terms to be counted are those with \(\mathfrak d(\gamma)\ne0\) and those
with \(\mathfrak d(\gamma)=0\), \(\gamma I\ne I\), and matching rational
unipotent characters. The latter matching condition means equality of
the character values of rational unipotent elements satisfying
\(u\gamma=\gamma u'\). This set is stable under the rational unipotent
actions on both sides.
\leansrc{New/Spanning/Unfolding/TermStructure}{The remaining terms and rational character matching}

The lattice \(L_\gamma\) of a near term \(\gamma\) consists of the vectors
in \(X_0\) in the lattice of class \(c\) that can be made orthogonal to the
frame of \(\gamma\) by adding \(M\) times an integral isotropic vector.
For two representatives \(\gamma,\gamma'\) over \((c,s,c')\), take the
intersection of their correspondences and impose the vanishing of the
last frame component in both coordinates. Project to \(X_0\) by the first
coordinate and restrict \(L_\gamma\cap L_{\gamma'}\) to this subspace.
Denote the resulting lattice by \(L_{12}\), its corank by \(l\), and its
discriminant by \(d\). The discriminant is the square root of the generalized
index of the Gram determinant of a lattice basis. Each correspondence has
dimension \(k+1\). If their intersection has codimension \(i\) in each
correspondence, the two last-coordinate conditions give \(l\le i+1\).
When \(i=0\), vanishing of one \(\mathfrak d\) implies vanishing of both, and
then \(l=0\). The bounds on lengths and lattices for near terms give
\(l\le n^3\) and \(1\le d\le b^{n^3 f}\).
\leansrc{New/Spanning/RemainingTerms/Gain}{Rank and discriminant bounds for the common lattice}

For a term of cell rank \(j\), define \(G_\gamma\) by taking the larger of
\(N^f\) times the normalized norm of the leading minor and the norm of
\(\mathfrak d(\gamma)\), multiplying by \(M^{2\min(\varrho,j)f}\), and dividing
by the product of the two indices of the isotropic row lattices.
For \(1\le z\le\operatorname{Meet}(s)\), set \(\zeta=z\) when \(i=0\) and
\(\mathfrak d(\gamma)\ne0\), and set \(\zeta=1\) otherwise. The discriminant
comparison for the two lattices gives
\begin{equation}\label{eq:span-pair-gain}
 (d\zeta)^2\le (2/M)^{2if}G_\gamma G_{\gamma'}.
\end{equation}
The factor for \(\zeta=z\) comes from a determinant. Over a residue field
where the frames meet, express a common vector in both frames and take
the largest index at which its two coefficients are not both zero.
Character matching allows the same pair of coefficients to occur at
the next index in a vector common to the two flags. Indeed, all linear
forms killing the columns before the next index in both frames take the
same value on the next columns weighted by these two coefficients.
Iteration gives a relation between the columns of both frames with at
least one nonzero last coefficient.
\leansrc{New/Spanning/Unfolding/Climbing}{Passing the coefficients of a common vector up the flags}
Pair this relation with the two integral lifts orthogonal to the earlier
part of the frames. The two last columns then have a nonzero kernel.
The resulting determinant \(\delta\) vanishes at every place counted in
\(\operatorname{Meet}(s)\). The quotient
\(\widehat{\mathcal O}_E/\delta\widehat{\mathcal O}_E\) maps onto the product
of the residue fields at these places, so the generalized index of
\(\delta\) is at least \(\operatorname{Meet}(s)\).
\leansrc{New/Spanning/RemainingTerms/Cut}{The determinant at places where the frames meet}

To obtain the comparison, split the common part of the correspondences
into the subspace cut out by the two last-coordinate conditions and its
complement. Then take the orthogonal complement of the common part in
each correspondence. Write \(d_\gamma^0\) for the discriminant of the
part of \(L_\gamma\) cut out by the last-coordinate conditions,
\(a_\gamma\) for the torsion index from the isotropic rows, and
\(\Delta_\gamma\) for the squared discriminant of the lattice outside the
common part. The determinant factor and the factorization of Gram
determinants give
\((a_\gamma d_\gamma^0\zeta)^2\Delta_\gamma\le G_\gamma^2\).
The two outer complements pair nondegenerately. Since the reductions
agree, the pairing defined by the difference of the Hermitian forms in
the two coordinates is integral after division by \(M\). Estimating its
determinant by the two Gram determinants gives
\(1\le(2/M)^{4if}\Delta_\gamma\Delta_{\gamma'}\).
Imposing the conditions of \(L_{\gamma'}\) on the common part of
\(L_\gamma\) increases its index by at most \(a_{\gamma'}\). Interchanging
the two terms gives the other comparison. Multiply
\(d\le d_\gamma^0a_{\gamma'}\) and \(d\le d_{\gamma'}^0a_\gamma\), then
eliminate the complementary discriminants to obtain
\eqref{eq:span-pair-gain}.
\leansrc{New/Spanning/RemainingTerms/Gain}{Discriminant comparison for two correspondences}

Fix the corank \(l\). A term of a given label in \(\bar R_0\) is
determined by its frame modulo the level-\(M\) unipotent moves of class
\(c\). To count frames with \(L\subset L_\gamma\), first fix the cell rank
\(j\) and the matrix \(B\) of frame products. Let \(B_j\) have the
first \(j\) rows of \(B\) and the remaining rows of the identity.
Its determinant \(\delta_j\) is the nonzero leading minor of order \(j\). Write \([a]\) for the index
\(|\widehat{\mathcal O}_E/a\widehat{\mathcal O}_E|\) of an integral finite idele \(a\).

Write a basis of \(L\) as \(Y=yA_L\), with \(k-l\) independent rational
columns in \(y\) and \(A_L\) invertible over the finite adeles.
The hypothesis \(L\subset L_\gamma\) gives a lift \(w\) satisfying
\[
 X_\gamma^*y+M B_{\mathrm{iso}}^*w=0,\qquad
 wA_L\text{ is integral},
\]
where \(X_\gamma\) is the matrix of \(X_0\)-components of the frame and
\(B_{\mathrm{iso}}\) consists of the isotropic rows of \(B\).
Here \({}^*\) denotes conjugate transpose. Reduce each row of \(wA_L\)
modulo the row lattice of \(Y^*Y\). That lattice has index
\(\operatorname{disc}(L)^2\), giving at most
\(\operatorname{disc}(L)^{2\varrho}\) classes for all \(\varrho\) rows.
Within one class a unipotent move makes the right-hand side of the
orthogonality condition for \(X_\gamma^*y\) the same for all frames.

The \(n\) columns of a difference of the matrices \(X_\gamma\) now lie
in the rank-\(l\) lattice orthogonal to \(L\). Reduction modulo
\(M\delta_j\) gives \([M\delta_j]^{ln}\) classes. Within one such class
unipotent moves make the \(X_0\)-components coincide. For the remaining
isotropic components, first fix the integral matrix of adjugate
entries recording the common vectors from stage \(j\) onwards. The
near condition bounds each entry by \(b^{182n^2}\) at infinity, giving
at most \(b^{366n^4 f}\) such matrices. Finally reduce the \(\varrho n\)
entries of the isotropic block modulo \(M\det(B_j B_j^*)\).
Frames in the same class coincide after a unipotent move.
\leansrc{New/Spanning/RemainingTerms/FrameParts}{The lift equation and unipotent moves}
\leansrc{New/Spanning/RemainingTerms/FramesOverLattice}{Congruence classes of frames over a fixed lattice}

For fixed \(j,B\), the number of these classes is at most
\[
 \operatorname{disc}(L)^{2\varrho}\,[M\delta_j]^{ln}
       b^{366n^4 f}[M\det(B_j B_j^*)]^{n\varrho}.
\]
Both modulus indices and \(\operatorname{disc}(L)^2\) are at most
\(b^{2n^3f}\). There are at most \(b^{2f}\) choices of \(j\) and
\(b^{270n^3f}\) choices of \(B\). Thus there are at most
\(b^{2000n^6(l+1)f}\) near representatives whose lattices contain \(L\).
A label can be chosen for each reduction with matching characters,
and the same label applies to both terms. There are consequently at
most \(b^{4000n^6(l+1)f}\le\tau\) ordered pairs.

To sum the weights \(d^{-k}\) for lattices of rank \(k-l\), average over
the cosets at maximal level, of which there are at least \((N^k/b)^f\).
In one unimodular lattice, extend a basis of the saturation of \(L\)
to a basis of the whole lattice. Write the basis of \(L\) as the
basis of its saturation multiplied by an integral matrix \(A_{\mathrm{sat}}\).
For the specified row to vanish on \(L\) modulo \(\pi_N\), its first
\(k-l\) values \(a\) must satisfy \(aA_{\mathrm{sat}}=\pi_N w\).
The class of \(w\) modulo the row lattice of \(A_{\mathrm{sat}}\) determines
\(a\) modulo \(\pi_N\), giving at most \([\det A_{\mathrm{sat}}]\) choices.
The other \(l\) values contribute
\(\bigl|\widehat{\mathcal O}_E/\pi_N\widehat{\mathcal O}_E\bigr|^l=N^{lf}\).
Distinct cosets have distinct rows modulo \(\pi_N\), so at most
\([\det A_{\mathrm{sat}}]N^{lf}\) cosets contribute for this lattice. Summed over the cosets, the weighted
lattice count with \(d\le D_\ast=b^{n^3f}\) is at most
\[
  N^{lf}D_\ast^{2\varrho+6}
       (200fk\,\operatorname{rd}_E)^{f(k+1)l}2^k
  \le b^{-f}\tau^{l+1},
\]
where \(\operatorname{rd}_E\) is the root discriminant of \(E\). Dividing by
the number of cosets bounds the average weight for one class by
\(N^{-kf}\tau^{l+1}\). Including all \(n_X\) classes, the \(|\bar R_0|\)
labels, and the multiplicity of pairs for each lattice gives
\begin{equation}\label{eq:span-pair-count}
 \sum_{\substack{\text{pairs of near representatives}\\
                            \operatorname{corank} L_{12}=l}}d^{-k}
 \le \tau^{l+2}n_p .
\end{equation}
\leansrc{New/Spanning/RemainingTerms/LatticesOfCorank}{Averaging lattice counts over cosets}
\leansrc{New/Spanning/RemainingTerms/Count}{Counting pairs weighted by their common lattice}

Schur's estimate for the orbit sum gives the bound
\[
 \sigma_\gamma=\tau^2\left(\frac{D_j}{n_pG_\gamma^k}\right)^{1/2},
 \qquad
 D_j=\begin{cases}
       D_{\mathrm{main}} & j=\varrho+1,\\
       \tau^{-16}D_{\mathrm{id}} & j\le\varrho
     \end{cases}
\]
for each near representative. More precisely, for a subset of the
remaining terms stable under both rational unipotent actions, the norm
of its truncated
amplitude divided by \(\nu\) is at most \(\sum_\gamma\sigma_\gamma\) over
the chosen near representatives. The case \(j=\varrho+1\) is precisely
\(\mathfrak d(\gamma)\ne0\).
\leansrc{New/Spanning/RemainingTerms/Size}{Amplitude bound for the double-orbit sum}
If \((d\eta)^2\le G_\gamma G_{\gamma'}\) and \(D\ge D_j,D_{j'}\), then
\[
 \sigma_\gamma\sigma_{\gamma'}
 \le \frac{\tau^4D}{n_p\eta^k}\,d^{-k}.
\]
Thus if \(\tau^{l+6}D/\eta^k\le C\) for each corank \(l\),
\eqref{eq:span-pair-count} bounds the sum of these weights at each corank
by \(C\). There are at most \(n^3+1\le\tau\) coranks, giving the bound \(\tau C\).

For equal correspondences with nonzero discriminants, \(l\le1\), and
\eqref{eq:span-pair-gain} allows \(\eta=z\) and \(D=D_{\mathrm{main}}\).
Taking \(C=\tau^7D_{\mathrm{main}}/z^k\) gives
\(\tau^8D_{\mathrm{main}}/z^k=(b^{2f}/z)^kD_{\mathrm{main}}\).
Each pair determines its reduction and its correspondence, so this
proves \eqref{eq:span-intersect}.

For equal correspondences with vanishing discriminants, take
\(l=0\), \(\eta=1\), and \(D=\tau^{-16}D_{\mathrm{id}}\).
Here \(C=\tau^{-10}D_{\mathrm{id}}\) suffices. For distinct correspondences,
\(i\ge1\) and \(l\le i+1\). Take \(\eta=(M/2)^{if}\) and
\(D=D_{\mathrm{main}}\). Since \(\rho_\ast=\tau(2/M)^{kf}\le1\) and
\(D_{\mathrm{main}}\le(M/b^6)^{kf}D_{\mathrm{id}}\), we have
\[
 \frac{\tau^{l+6}D_{\mathrm{main}}}{\eta^k}
 \le \tau^7\rho_\ast^iD_{\mathrm{main}}
 \le \tau^8(2/b^6)^{kf}D_{\mathrm{id}}
 \le \tau^{-10}D_{\mathrm{id}} .
\]
The last inequality uses \(2\le b\) and \(\tau=b^{kf/4}\).
Both cases therefore fit the same bound \(C=\tau^{-10}D_{\mathrm{id}}\),
and their total is at most \(\tau^{-9}D_{\mathrm{id}}\le b^{-kf}D_{\mathrm{id}}\).
\leansrc{New/Spanning/RemainingTerms/Sum}{Squares for one correspondence and the cross terms}

To apply the count, split each block into one piece per correspondence
with nonzero discriminant and a single piece of all terms with vanishing
discriminant. Square the sum of the norms of the pieces. Summing the
diagonal terms for the pieces of nonzero discriminant gives
\(\mathcal L(\Omega)\). Use the products
\(\sigma_\gamma\sigma_{\gamma'}\) for the other terms. A pair belongs to
only one reduction, so the same bound applies after summing over the blocks.

For the truncation error, the Bessel tail bound gives
\(q=\nu/(2b^{k^3 f})\) as an upper bound on the sum of the norms of the
amplitude differences in each row \((c,s)\). The sum of squares in the
row is at most \(q^2\). The number of rows and the inequalities
\(|\bar G|b^{40}\le(2b^{k^3 f})^2\) and \(n_X\le D_{\mathrm{id}}\) show that
the sum over all rows divided by \(\nu^2\) is at most
\(b^{-40}D_{\mathrm{id}}\). This estimate compares the full and truncated
functions on the same set of terms.
\leansrc{New/Spanning/RemainingTerms/Size}{Row estimate for the truncation error}
On the allowed set \(\Omega\), the full amplitude for \(E_s\) equals the
sum over the remaining terms. Outside the Bessel subgroup and the
remaining terms, those failing the rational character matching condition
give a vanishing sum over a stable set. Those with matching rational
characters have a norm-one coefficient pair from the first stage, and
their reductions lie in the excluded scalar cosets.
\leansrc{New/Spanning/Unfolding/SquareExpansion}{The amplitude outside the Bessel subgroup}
The triangle inequality for the square root of the sum of squares and
\(kf\ge40\) give
\[
 \Bigl(\sum_{s\in\Omega}\|E_s\|_{\mathrm{HS}}^2\Bigr)^{1/2}
 \le \mathcal L(\Omega)^{1/2}
       +(b^{-kf/2}+b^{-20})D_{\mathrm{id}}^{1/2}
 \le \mathcal L(\Omega)^{1/2}+b^{-15}D_{\mathrm{id}}^{1/2}.
\]
This proves \eqref{eq:span-cutoff}.

\subsection{The second moment on open reductions}\label{sec:span-open-proof}

\paragraph{Local mass under fixed-vector conditions}
Choose a conjugation-stable finite set of finite places \(\mathcal T\)
containing the support of the level \(M\). Write \(K_{\mathcal T}\) for the
integral unitary group over the product of completions and
\(K_{\mathcal T}(g)\) for its principal congruence subgroup of local level
\(g\). For \(K_{\mathcal T}(g)\subset Z\subset K_{\mathcal T}\) and Haar
measure \(\mu\), put
\[
 \mathfrak M(Z)=\int
    \left|\int_Z \chi_g(yxy^{-1})\,d\mu(y)\right|^2d\mu(x),
 \qquad
 m_Z=\frac{\mu(Z)\mu(K_{\mathcal T}(g))}{u_g^2}.
\]
Here \(\chi_g\) is the unipotent character extended to the product of its
subgroup with \(K_{\mathcal T}(g)\) and by zero elsewhere, and \(u_g\) is
the associated unipotent measure constant. For \(x\in K_{\mathcal T}(g)\),
\(\chi_g(yxy^{-1})=1\), so the contribution from this subgroup is exactly
\(\mu(Z)^2\mu(K_{\mathcal T}(g))\).

Suppose the local lattice data have the following property for \(Q\ge2\).
An invertible integral element whose inverse is nonintegral has modulus
at most \(Q^{-1}\). The estimate outside the congruence subgroup divides
the integral into ranges of the index of the denominator ideal. For
\(a+10^8n^4\le k\), the integral of the zeroth-stage majorant there is at most
\(m_Z2^{10^9n^4}\min(1,Q^{-a})\). If no coordinate at a position \(x\) outside the congruence subgroup
had depth at least \(Q\), one could choose integral units among the
stable frame products. This would make \((x-1)/g\) integral, a contradiction.
Thus the bound for a deep vector applies. The factor in the integral
estimate for the number of coordinates is at most \(2^{99\cdot13n}\).
The sum over dyadic ranges of the denominator index incurs at most
\(2^{4\cdot10^8n^4}\). Their product is at most \(2^{10^9n^4}\).
Taking \(a=t+10^9n^4\) absorbs the extra factor by
\(2/Q\le1\). Multiplying the majorant by \(u_g^2\) and adding the principal
congruence contribution gives, for \(t+2\cdot10^9n^4\le k\),
\begin{equation}\label{eq:span-local-mass}
 \mathfrak M(Z)
 \le\mu(Z)\mu(K_{\mathcal T}(g))\,\bigl(\mu(Z)+Q^{-t}\bigr).
\end{equation}
\leansrc{New/Spanning/Frames/FiniteShells}{The local integral estimated by denominator ideals}

On the factor for a single \(w\in A\) and its conjugate, denote the
integral group by \(K_w\) and the inverse image of the fixed-vector
condition by \(Z_w\). Reduction onto \(\mathrm{GL}_r(\mathbf F_{p_0})\) is
surjective, and the conjugate factor acts by the inverse adjoint.
Transitivity on nonzero vectors in either action gives
\([K_w:Z_w]=p_0^r-1\). Put \(B=b^9\), so \(2\le B\le p_0\).
The additive volume of nonunits in the product over the selected places
is at most \(B^{-1}\). Each residue field has cardinality at least
\(b^{10}\) and contributes at most its reciprocal, with at most \(b\) places.
The same bound applies to the modulus of an element whose inverse is nonintegral.
The admissibility conditions give \(t_0=k-10^{10}n^4\) with
\[
 t_0+1+n+2\cdot10^9n^4\le k,\qquad
 6k\le9t_0,\qquad t_0+1<r.
\]
The index formula gives \(\mu(Z_w)\le B^{-t_0-1}\mu(K_w)\).
In the normalization of the local integral,
\(\mu(K_w)\ge2^{-n}\ge B^{-n}\). Apply \eqref{eq:span-local-mass} with
\(Q=B\) and \(t=t_0+1+n\) to obtain
\[
 \mathfrak M(Z_w)\le
 \mu(Z_w)\mu(K_w(g))
       \bigl(\mu(Z_w)+B^{-t_0-1}\mu(K_w)\bigr).
\]
Multiply by \(\mu(K_w)b^{6k}\) and use \(b^{6k}\le B^{t_0}\) and \(2/B\le1\).
The contribution to \(\mathfrak M(K_w)\) from \(K_w(g)\) alone is
\(\mu(K_w)^2\mu(K_w(g))\), so
\[
 \mathfrak M(Z_w)\mu(K_w)b^{6k}
 \le\mathfrak M(K_w)\mu(Z_w).
\]
The groups and characters factor over distinct conjugate pairs. Multiply
the mass integrals and volumes, taking the full integral group on the
remaining factors. This gives the same inequality with
\(K_w,Z_w,b^{6k}\) replaced by \(K_{\mathcal T},Z,b^{6k|A|}\).
The ratio comparison is invariant under a common rescaling of the
product Haar measure.
\leansrc{New/Spanning/Frames/FixedLines}{Mass ratio for each fixed-vector condition}

For \(Z=K_{\mathcal T}\), apply \eqref{eq:span-local-mass} with \(Q=2\)
and \(t=n\). The inequalities \(\mu(K_{\mathcal T})\le1+2^{-n}\) and
\(2^{-\varrho}\le b^{-6}\) give
\(\mathfrak M(K_{\mathcal T})\le(1+b^{-6})\mu(K_{\mathcal T})\mu(K_{\mathcal T}(g))\).
The zeroth-stage orbital form is a squared correlation averaged against
the character, and its unnormalized absolute value times \(u_g^2\) is
at most \(\mathfrak M(Z)\). Dividing by \(m_Z\) bounds the normalized
zeroth-stage finite-place kernel by
\begin{equation}\label{eq:span-local-zero}
 m_A=(1+b^{-6})b^{-6k|A|}.
\end{equation}
\leansrc{New/Spanning/MainTerm/FiniteSieveSystem}{Zeroth-stage finite-place density}

\paragraph{Siegel's estimate at a fixed last entry}
Take finitely many nonzero amplitude blocks \((c,s,c',\Lambda)\) over
\(\Omega_A\). All proper leading minors are nonzero at the rational
point, so the big-cell decomposition names each block by its last
entry \(t\in E^\times\), a representative \(\sigma\) fixing the preceding
hyperbolic basis vectors, and an integral lift \(n_s\).
They satisfy \(|\Phi_v(t)|^2>N^2\) at every infinite place \(v\), and
\((c,c',\operatorname{red}(n_s),\sigma)\) determines the block.
With \(\sigma\) fixed, the remaining leading entries are in a finite set
\(T_0\subset E^\times\). The block amplitude can be written
\(\mathcal A=\sum_{\mathbf p\in T_0^\varrho} a(\mathbf p)\), where
\(a(\mathbf p)\) is the truncated unipotent double-coset amplitude of the
Weyl representative \(w(\mathbf p,t)\sigma\). Enlarge \(\mathcal T\) so
that the elements of \(T_0\) and their inverses, \(M^{-1}\), \(N^{-1}\),
the lattice basis and its inverse, and the class translates of each
\(\sigma\) are integral outside \(\mathcal T\).
\leansrc{New/Spanning/MainTerm/EnergyEstimate}{Big-cell amplitude sums and reduction to finite sets}

At an infinite place \(v\), let
\(\lambda_{v,j}=\kappa_v^{2(\varrho-j)}/N\) be the scale of entry \(j\), where
\(\kappa_v=M\Phi_v(\varpi)/(2\pi\mathrm i)\) for the generator \(\varpi\) of the
different. Let \(R_t(\mathbf p,\mathbf q)\) be the product of the top
kernels at the infinite places at last entry \(t\). These are Gram kernels
of orbital integrals with weights given by the coordinate scales and the
densities of the preceding entries. The last-entry finite-place density
\(F_t(\mathbf p,\mathbf q)\) is an integral over a quadric of the Gram
pairings on fibers with a common last vector.

For a fixed \(t\), the \(X_0\)-component \(y\) of the image of the last
frame vector under \(\sigma\) satisfies \(N\langle y,y\rangle=t\bar t-N^2\).
Set
\[
 G_t=\left(\prod_v(|\Phi_v(t)|^2-N^2)\right)^{k-1}>0 .
\]
Write \(d_{B,v}\) for the number of Bessel coordinates at \(v\) and \(d_B\)
for their product. With the big-cell volume constant
\(c_G=\operatorname{vol}(U(k))^2/
(\operatorname{vol}(U(r))\operatorname{vol}(U(k-1)))\), the weight of the
archimedean orbital integral is
\[
 \rho_t(\mathbf p)=
 \prod_v \frac{\sqrt{c_G}}{d_{B,v}}
 \left(\frac{|\Phi_v(t)|^2}{N^2}-1\right)^{(k-1)/2}
 \prod_{j<\varrho}|\lambda_{v,j}\Phi_v(p_j)|^{k-1+2(\varrho-j)} .
\]
Let \(h_t(\mathbf p)\) be the corresponding finite-place density of
the head. Since \(p_j\) is a unit outside \(\mathcal T\), the product
formula identifies \(\prod_v|\Phi_v(p_j)|\) with the inverse square root
of the modulus of \(p_j\) at \(\mathcal T\). This gives
\(\rho_t(\mathbf p)h_t(\mathbf p)=r_0\sqrt{G_t}\) with
\[
 r_0=\frac{1}{d_B}\prod_v
       \left(\sqrt{c_G}\prod_{j<n}\lambda_{v,j}^{k-1+2(\varrho-j)}\right).
\]
Put \(D=|\bar R_0|M^{f\varrho(2k+2\varrho+1)}\) and \(C_S=D\,n_Xc_S\),
where \(c_S\) is the constant in the Siegel sum on the sphere.
We first sum the blocks with a common point \(x=(c,y)\).
By \(\Psi_t(x)\) we mean evaluation at the last local vector obtained
from \(y\) after transporting it by the inverse of the class representative.
Put \(\mathcal K=K_{\mathcal T}(g)\) and \(\mathcal C=Z/\mathcal K\).
The scaled orbital integral for a head \(\mathbf p\) has coordinates
in the finite set \(\mathcal C^2\). For a position \(\ell\) and a
coordinate \((a,a')\in\mathcal C^2\), denote the vector of these
coordinates, indexed by the heads, by \(v_{\ell,a,a'}\).
The sum of the finite-place pairings weighted by \(R_t\) is
\[
 Q_t(\ell)=\sum_{a,a'\in\mathcal C} v_{\ell,a,a'}^*R_t v_{\ell,a,a'}\ge0 .
\]
For fixed \(c,c',\sigma\), the block with reduction \(s\) occurs in
the coordinates where \(a'\) is the coset of \(n_s^{-1}a\).
Distinct reductions occupy disjoint sets of coordinate pairs.
Unfolding the orbital integral identifies each such summand with
\(\mu(\mathcal K)^2G_t r_0^2\|\mathcal A\|^2\). Summing over \(a\) gives
a factor \(|\mathcal C|\), and the sum over all the blocks is bounded by
\(Q_t(\ell)\).

Write \(H_{\mathrm{fix}}\) for the local stabilizer of the last frame
vector in the subgroup fixing the hyperbolic pairs. The positions
\(\ell\) for \((c',\sigma)\) with the same \(x\) lie in a single orbit
under left multiplication by \(H_{\mathrm{fix}}\). They lie in distinct
congruence cosets. Indeed, congruence at \(\mathcal T\) together with
integrality away from \(\mathcal T\) identifies both the source class
and the rational representative. The integrand \(Q_t\) is invariant
under left multiplication by \(H_{\mathrm{fix}}\cap\mathcal K\).
Multiplying the sum by the measure of a congruence coset thus bounds it
by the integral on the fiber. For the Haar measure \(\mu_H\) on
\(H_{\mathrm{fix}}\), the mass comparison gives
\(D\mu_H(H_{\mathrm{fix}}\cap\mathcal K)\mu(\mathcal K)^2|\mathcal C|\ge m_Z\).
The step function \(\Psi_t\) extending the finite-place fiber pairing
therefore satisfies
\[
 G_t(m_Z/D)\sum_{(c,y)=x} r_0^2\|\mathcal A\|^2\le\Psi_t(x).
\]
\leansrc{New/Spanning/Frames/TopFloor}{Finite quotient coordinates and the fiber pairing}

Put \(\beta_t=(t\bar t-N^2)/N\). Apply Siegel's formula for one vector in
the definite space of rank \(k\ge3\). The test function imposes
integrality away from \(\mathcal T\) and pulls back \(\Psi_t\) along the
last-vector map at \(\mathcal T\). It is compactly supported, locally
constant, and invariant under the source level. The norm \(\beta_t\)
is totally positive. The class sum is the finite adelic norm density
times
\[
 n_X|\operatorname{disc}(E)|^{1/4}
       \prod_v\frac{\pi^k(\Phi_v(\beta_t))^{k-1}}{(k-1)!}.
\]
\leansrc{Lib/SiegelWeil/OneVector/SiegelClassSum}{Siegel's formula at one norm}

The norm density integrates the test function times
\(\psi_f(u(w^*w-\beta_t)/2)\) over the finite vectors \(w\) and
self-adjoint finite adelic frequencies \(u\). Splitting the vectors
and frequencies by their components at \(\mathcal T\) factors the
density into the sphere density there and an integral of Gauss sums
from the integrality condition outside \(\mathcal T\). Write
\[
 F_{u,\Gamma}(w)=\psi_f\!\left(u(w^*w-\Gamma)/2\right),
 \qquad w\in\widehat{\mathcal O}_E^{\,k}.
\]
Here \(\Gamma\) is the shift from the prescribed norm that remains in
the Gauss sum, and the bound is uniform in \(\Gamma\).
The square of the absolute value of \(\int F_{u,\Gamma}(w)\,dw\)
is the integral of \(\overline{F_{u,\Gamma}(w)}F_{u,\Gamma}(w+h)\)
over \(w,h\). Translating coordinate \(j\) of \(w\) by \(\bar a\)
multiplies the correlation by \(\psi_f(u h_j a)\).
Consequently the integral over \(w\) vanishes unless every \(h_j\)
lies in the conductor ideal
\[
 I_u=\{h\in\widehat{\mathcal O}_E:
            \psi_f(uha)=1\text{ for all }a\in\widehat{\mathcal O}_E\}.
\]
The integrand has modulus \(1\), so the remaining correlation is
bounded by the product of the volumes. With
\(q_u=|\widehat{\mathcal O}_E/I_u|\), it follows that
\[
 \left|\int F_{u,\Gamma}(w)\,dw\right|
 \le \operatorname{vol}(\widehat{\mathcal O}_E)^k q_u^{-k/2}.
\]
For a fixed conductor \(I\), the frequency classes modulo the
self-adjoint codifferent are distinguished by the characters of
the finite additive group \(\widehat{\mathcal O}_E/I\). Each class has
measure \(1\), and there are at most \([\widehat{\mathcal O}_E:I]\) of them.
With \(\mu_0=3\varrho+4\) and \(k\ge2(\mu_0+1)\), integration over the
frequencies is therefore bounded by
\(\operatorname{vol}(\widehat{\mathcal O}_E)^k
\sum_I[\widehat{\mathcal O}_E:I]^{-\mu_0}\).
The sum over finite-index ideals is at most
\(1+2^{-\varrho}\le1+b^{-6}\), by comparison with the Dedekind zeta
Euler product using \([E:\mathbf Q]\le2^{2\varrho}\).
\leansrc{Lib/SiegelWeil/Mass}{Character orthogonality and the Gauss-sum bound by conductor}

Finite-place Fourier inversion bounds the density at \(\mathcal T\)
by the integral of the test function on the sphere
\(\langle y,y\rangle=\beta_t\). The density \(F_t\) divides this integral by
the modulus \(N^{-f}\) of \(\pi_N\). Expressing the integral in terms of
\(F_t\) multiplies the archimedean norm factor \(N^{-f(k-1)}\) by \(N^{-f}\),
giving
\[
 c_S=\left(\frac{\pi^k}{(k-1)!N^k}\right)^f
       |\operatorname{disc}(E)|^{1/4}
       \operatorname{vol}(\widehat{\mathcal O}_E)^k.
\]
Summing over the distinct points \(x\) gives
\[
 \sum_x\Psi_t(x)\le(1+b^{-6})n_X c_S G_t\,
             \operatorname{Re}\sum_{\mathbf p,\mathbf q}R_t(\mathbf p,\mathbf q)
                                                       F_t(\mathbf p,\mathbf q).
\]
Cancel \(G_t\), divide by \(m_Z\), and sum over \(t\) to obtain
\begin{equation}\label{eq:span-siegel-sum}
 r_0^2\sum_{\mathrm{blocks}}\|\mathcal A\|^2
 \le (1+b^{-6})C_S\sum_t\operatorname{Re}\sum_{\mathbf p,\mathbf q}
                         R_t(\mathbf p,\mathbf q)\frac{F_t(\mathbf p,\mathbf q)}{m_Z}.
\end{equation}
\leansrc{New/Spanning/MainTerm/SiegelTopBound}{Siegel's estimate for the class sum and its local densities}

\paragraph{Removing the entries by Poisson summation}
Let \(J_i\) be the product of the archimedean kernels and the normalized
finite-place kernel at stage \(i\). Let \(\mathcal S_{i+1}\) be the real part of
the sum of \(J_{i+1}\) over two heads \(\mathbf p,\mathbf q\in T_0^i\)
followed by a common entry \(t\in T_0\). Entries after the heads are filled
with \(1\), and each entry at infinity is multiplied by its scale
\(\lambda_{v,j}\). At stage zero \(\mathcal S_0=\operatorname{Re}J_0\) on the empty
heads. The sum for each common entry is nonnegative by the Gram property,
so the kernel sum on the right of \eqref{eq:span-siegel-sum} is at most \(\mathcal S_n\).

Apply the one-sided Poisson summation bound to entry \(i\).
The sum of the kernels with that entry held in common is a nonnegative
lower semicontinuous function. Multiply it by the indicator of the
integral adeles outside \(\mathcal T\). By the integral identity for the
orbital forms, the zero frequency gives the sum of \(J_i\) over all
pairs of heads, with multiplier
\[
  q_i=\frac{\operatorname{vol}(\widehat{\mathcal O}_E)}
                 {\prod_v\lambda_{v,i}^2}.
\]
The other frequencies are rational points integral outside \(\mathcal T\)
and scaled by \(\varpi^{-1}\). Let \(\epsilon\) bound the sum of the
products of the local Fourier decay factors at the nonzero frequencies,
and let \(W_i\) be the sum of the nonnegative weights of heads of length \(i\).
The contribution from those frequencies is at most \(q_i\epsilon W_i^2\).
In the sum of \(J_i\) produced by the zero frequency, the pairs with
equal last entries give \(\mathcal S_i\). For unequal last entries \(u,v\), the
finite-place frequency \(M^{-1}(u^{-1}-v^{-1})\) is nonzero and integral
outside \(\mathcal T\). The same decay estimate bounds these pairs by
\(\epsilon W_i^2\), giving
\[
 \mathcal S_{i+1}\le q_i(\mathcal S_i+2\epsilon W_i^2).
\]

For the frequency bound put \(K_0=k-n^7\).
At infinity the decay function is \(d_{K_0}(a)=a^{K_0}\min(1,a^6)\).
The factor at \(v\) and its volume ratio are
\[
 c_{\mathrm{vol}}b^{2000n^9}d_{K_0}(|\xi_v|^{-1}),
 \qquad
 c_{\mathrm{vol}}=\pi^{(2\varrho+1)k}
            \frac{\operatorname{vol}(U(k))}{\operatorname{vol}(U(r))}.
\]
At the finite places the factor is the generalized index of the frequency
to power \(K_0/2\), divided by its denominator index cubed, times \(b^{n^9 f}\).
The frequencies associated to a nonzero \(\mathcal T\)-integer \(y\)
have the form \(\xi_v=(c_v/\lambda_{v,j})\Phi_v(y)\) at infinity.
In both error terms the scales satisfy \(|c_v|^{-1}\le|\kappa_v|/M\), so
\[
 |c_v/\lambda_{v,j}|^{-1}\le
 \theta_v:=|\kappa_v|^{2\varrho+1}/(MN)
 \le B_F:=b^{23(2\varrho+1)}.
\]
Write \(H(y)=\operatorname{Norm}(\mathfrak a_y)\prod_v\max(1,|\Phi_v(y)|)^2\)
for the height of \(y\), where \(\mathfrak a_y\) is its denominator ideal.
Since \(y\) is integral outside \(\mathcal T\), the product formula bounds
the square root of the finite-place generalized index by
\(\prod_v|\Phi_v(y)|\). This upper bound cancels the factors
\(|\Phi_v(y)|^{-K_0}\) at infinity.
Since \(\min(1,|\xi_v|^{-1})\le B_F/\max(1,|\Phi_v(y)|)\) and the
finite-place denominator index is at least \(\operatorname{Norm}(\mathfrak a_y)\),
the contribution of one frequency, with the powers of \(b\) and volume
ratios factored out, is at most \(H(y)^{-3}\prod_v(\theta_v^{K_0}B_F^6)\).
The inverse cubes of the heights sum to at most
\((64\cdot38\exp(4))^f\le b^{3f}\).
Together with \((bB_F^2)^3\le b^{n^9}\), this gives
\[
 \epsilon\le b^{2002n^9f}\mathcal D_\ast,\qquad
 \mathcal D_\ast=\prod_v c_{\mathrm{vol}}\theta_v^{K_0}.
\]
\leansrc{New/Spanning/MainTerm/GlobalSieve}{Summing frequencies using the product formula and height}

For the head weight, each prefix product contributes its finite-place
generalized index to power \(n^6/2\). At each infinite place, it
contributes the larger of \(1\) and the inverse modulus of the scaled
prefix product, raised to \(n^6\). The weight is zero if a prefix product
is nonintegral or if its scaled modulus exceeds \(b^{60n^2}\).
A head of nonzero weight is determined by its sequence of prefix
products, which are algebraic integers of modulus at most \(b^{120n^2}\)
in every embedding. The lattice count bounds the number of such heads
by \((1+2b^{120n^2})^{2fi}\). By the product formula each weight is at
most \(b^{120n^8fi}\), giving \(W_i\le b^{130n^9f}\) for \(i\le\varrho\).
Also \(\lambda_{v,j}\le b^{46\varrho}\) and
\(\operatorname{vol}(\widehat{\mathcal O}_E)=(2/\operatorname{rd}_E)^f\)
give \(\prod_{j<i}q_j\ge b^{-100n^2f}\). Hence
\[
 W_i^2\le b^{260n^9f+100n^2f}\prod_{j<i}q_j
       \le b^{400n^9f}\prod_{j<i}q_j .
\]
\leansrc{New/Spanning/MainTerm/GlobalSieve}{Counting prefix products and bounding head weights}
Write \(\deg_v\) for the formal degree of the Bessel system at \(v\).
Using \(\mathcal S_0\le\prod_v((1+b^{-10})/\deg_v)\,m_A\) and
\(2n\le b^{n^9f}\), iterate the recurrence \(n\) times to obtain
\begin{equation}\label{eq:span-sieve-sum}
 \mathcal S_n\le \left(\prod_{j<n}q_j\right)
  \left[
    \left(\prod_v\frac{1+b^{-10}}{\deg_v}\right)m_A
    + b^{2403n^9f}\mathcal D_\ast
  \right].
\end{equation}

\paragraph{Comparison of constants}
Let \(\Lambda_\ast\) be the product of the entry scales to powers
\(k+2(\varrho-j)\). The formulas for \(r_0\) and the Poisson multipliers
\(q_j\) give
\[
 V_\ast=\frac{\prod_{j<n}q_j}{r_0^2}
 =\frac{\operatorname{vol}(\widehat{\mathcal O}_E)^n}{c_G^f}
    \left(\frac{d_B}{\Lambda_\ast}\right)^2 .
\]
Substituting the covolume formula for \(\nu\) into the comparison of
Siegel constants gives
\[
 \begin{split}
 C_S V_\ast
  &\le(1+b^{-9})D_{\mathrm{main}}\deg_B\nu^2,\\
 C_S V_\ast\mathcal D_\ast
  &\le b^{n^9f}d_BD_{\mathrm{id}}\nu^2 .
 \end{split}
\]
The first inequality substitutes the class count and volumes into the
main term. For the second, put \(s_d=\sqrt{\operatorname{rd}_E}\),
\(Q_d=Ms_d^2/(2\pi)\), and \(x_d=Q_d^{2\varrho+1}/(MN)\).
Use \(\prod_v|\kappa_v|=Q_d^f\) and cancel the \(k\)-th power in
\(x_d^{k-n^7}\) against the powers in the Siegel constant. The ratio
\(C_SV_\ast\mathcal D_\ast/(d_BD_{\mathrm{id}}\nu^2)\) is the \(f\)-th
power of
\[
 \frac{(2/s_d^2)^{n+2\varrho^2}(s_d/s_d^{2\varrho})}
      {2M^{\varrho(2\varrho+1)}}\,x_d^{-n^7}.
\]
Since \(2/s_d^2\le b\), \(s_d/s_d^{2\varrho}\le1\), and
\(x_d^{-1}\le MN\le b^{23+60n}\), this factor is at most \(b^{n^9}\).
\leansrc{New/Spanning/MainTerm/EnergyEstimate}{Comparing the covolume and decay constants}

Put \(Y=b^{-6k|A|}D_{\mathrm{main}}\nu^2\). The admissibility inequalities
\(d_B\le b^{2kf}\), \(b^{11kf}D_{\mathrm{id}}\le D_{\mathrm{main}}\),
\(|A|\le f\), and \(512000n^9\le k\) give
\(2404n^9f+2kf+10+6k|A|\le11kf\). Thus the error in
\eqref{eq:span-sieve-sum}, after accounting for its normalization, is
\(C_SV_\ast b^{2403n^9f}\mathcal D_\ast\le b^{-10}Y\).
For the main term use \((1+b^{-10})^f\le1+b^{-9}\) and
\(\prod_v\deg_v=\deg_B\). The former follows from \(f\le b/1000\), which
gives \((1+b^{-10})^f\le\exp(b^{-9}/1000)\le1+b^{-9}\). Combining \eqref{eq:span-siegel-sum},
\eqref{eq:span-local-zero}, and \eqref{eq:span-sieve-sum} yields
\[
 \sum_{\mathrm{blocks}}\|\mathcal A\|^2
 \le (1+b^{-6})\left[(1+b^{-9})^2(1+b^{-6})+b^{-10}\right]Y
 \le (1+b^{-4})Y .
\]
For the last inequality, \(b\ge1000\) implies
\(b^{-6},b^{-9},b^{-10}\le s_\ast=b^{-4}/100\le1\), so the coefficient
is at most \((1+s_\ast)^4+s_\ast(1+s_\ast)\le1+17s_\ast\).
Taking the supremum over the finite sums and dividing by \(\nu^2\)
proves \eqref{eq:span-open-input}.

An element of \(\bar R_0G_i\) agrees with an element of \(\bar R_0\) above
the other prime. The frames therefore meet at all \(f\) chosen places
above that prime and at their conjugates, so \(z\) may be taken to be
that prime to the power \(2f\). The set also avoids the cosets of
\(\bar R_0\) obtained from nontrivial roots of unity. Using \(p_j^2\geq M/b\)
for both primes and \(D_{\mathrm{main}}\le (M/b^6)^{kf}D_{\mathrm{id}}\), the
right-hand side of \eqref{eq:span-intersect} is at most
\(b^{-3kf}D_{\mathrm{id}}\). Under \(kf\geq 40\) and \(b\ge 1000\),
\eqref{eq:span-cutoff} gives
\begin{equation}\label{eq:span-bessel-near}
 \sum_{s\in\bar R_0G_i}\|E_s\|_{\mathrm{HS}}^2
 \leq b^{-12}D_{\mathrm{id}}.
\end{equation}
\leanref{secondMoment_compl_besselS_le}

\subsection{Second moments before and after projection}

For \(t\in\bar R_0\), \eqref{eq:span-equivariance} gives
\(\|E_t\|_{\mathrm{HS}}=\|P^*P-1\|_{\mathrm{HS}}\).
Restricting \eqref{eq:span-bessel-near} to \(\bar R_0\) thus yields
\begin{equation}\label{eq:span-gram-isometry}
 \|P^*P-1\|_{\mathrm{HS}}^2\leq b^{-12}D_X.
\end{equation}
For the orthogonal projection \(Q_U\) onto a subspace \(U\subset S_0\),
\[
 \|P Q_U\|_{\mathrm{HS}}^2
 = \dim U+\operatorname{tr}\bigl(Q_U(P^*P-1)\bigr)
 \geq \dim U-b^{-6}D_X.
\]
To estimate the loss from projecting to the target, factor
\(e_0=e_{p_0}e_{p_1}\) and let \(a_i(t)\) be the coefficient of \(t\) in the
complementary projection \(1-e_{p_i}\) in \(\mathbf C[G_i]\).
Counting the types of the finite general linear groups specified by
polynomials and using their degrees gives
\[
 \sum_{t\in G_i}|a_i(t)|^2=a_i(1)\leq \tfrac12 b^{-8}.
\]
The equality follows by comparing coefficients of the identity in the
square of the self-adjoint idempotent. \leanref{bad_types_rare}

Since \(1-e_{p_i}\) is an orthogonal projection,
\[
 \|(1-e_{p_i})P\|_{\mathrm{HS}}^2
  =\operatorname{tr}\sum_{t\in G_i} a_i(t)\bigl(r_\psi(t)+E_t\bigr).
\]
The trace of the Bessel term is \(D_X\) times the reduced unipotent
character. Since \(a_i\) is a class function, this character can be
replaced by its conjugacy average \(\psi_U^\natural\). The bound
\(\sum_{s\in\bar G}|\psi_U^\natural(s)|^2\le2\) and Cauchy--Schwarz
show that the absolute value of this contribution is at most \(b^{-4}D_X\).
\leanref{trace_besselTranslate_eq_mul_reducedChar}

For the remaining term, extend \(s\mapsto E_s\) by zero outside \(\bar R_0G_i\)
and apply \(f(s)\mapsto\sum_t a_i(t)f(st)\). As an orthogonal projection on
\(\ell^2(\bar G)\), this operation does not increase the sum of squares.
At \(s\in\bar R_0\) its value is \(r_\psi(s)\sum_t a_i(t)E_t\) by
\eqref{eq:span-equivariance}. Hence \eqref{eq:span-bessel-near}, after division by
\(|\bar R_0|\), gives \(\|\sum_t a_i(t)E_t\|_{\mathrm{HS}}^2\le b^{-12}D_X\).
Its trace has absolute value at most \(b^{-6}D_X\le b^{-4}D_X\).
Applying this at the two primes and using the orthogonal decomposition
\(1-e_{p_0}e_{p_1}=(1-e_{p_0})+e_{p_0}(1-e_{p_1})\) gives
\(\|(1-e_0)P\|_{\mathrm{HS}}^2\le4b^{-4}D_X\). Therefore
\begin{equation}\label{eq:span-near-isometry}
 \|e_0P Q_U\|_{\mathrm{HS}}^2\geq \dim U-5b^{-4}D_X
       \qquad(U\subset S_0).
\end{equation}
\leanref{isNearIsometry_piTensorElem_comp}
Taking \(U=S_0\) in \eqref{eq:span-near-isometry} gives \(h>0\).

We next estimate the second moment over the full group. Put
\(\beta(s)=\|P_s\|_{\mathrm{HS}}^2\) and \(X_t=\rho(t)PP^*\rho(t)^*\).
Cyclicity of the trace gives
\[
 \langle X_t,X_u\rangle_{\mathrm{HS}}=\beta(t^{-1}u).
\]
Thus \(\beta\) is of positive type, and \(\beta(s)\leq\beta(1)\).
By \eqref{eq:span-gram-isometry}, we also have
\(\beta(1)\le(1+b^{-6})^2D_X\le2D_X\).

Among the finitely many choices of \(A\) and of its specified vectors,
choose one maximizing
\[
 b^{6k|A|}\sum_{s\in\Gamma_A}\beta(s).
\]
Put \(Q=\sum_{s\in\Gamma_A}\beta(s)\). First suppose \(A\) includes all the
chosen places above \(p_0\), and put \(C=b^{-6kf}D_{\mathrm{main}}\).
The union of the scalar cosets of \(\bar R_0\) obtained from integral adeles
has at most \(M^f|\bar R_0|\) elements. Its contribution is at most
\(2 M^fD_{\mathrm{id}}\le C/2\), by the parameter comparison
\(8 M^f D_{\mathrm{id}}\le b^{-6kf}D_{\mathrm{main}}\).
For the remaining elements we can take \(z=p_0^f\).
Since \(p_0\ge b^{10}\), the right-hand side of \eqref{eq:span-intersect}
is at most \(C/8\). The square of the error term in
\eqref{eq:span-cutoff} is also at most \(C/8\), so the square of the sum
of the two square roots is at most \(C/2\). It follows that \(Q\le C\).
\leanref{secondMoment_fixer_univ_le}

Now suppose there is a chosen \(w\notin A\). By maximality, the sum of
\(\beta\) over the elements fixing any specified nonzero frame vector
at \(w\) or its conjugate is at most \(b^{-6k}Q\).
We use this to bound the mass of the singular minors at \(w\).
Let \(\pi\) be the matrix representation of \(\Gamma_A\) over the residue
field of order \(p_0\). Write the \(\ell\)-th leading principal minor as
\(\det Z_s\), where \(Z_s=A_1^{\mathsf T}\pi(s)A_0\). The columns of
\(A_0\) are the first \(\ell\) frame vectors, and the columns of \(A_1\)
represent the corresponding linear forms. Orthogonality for a nontrivial
additive character \(\psi\) gives
\[
 p_0 \sum_{\det Z_s=0}\beta(s)
 \le \sum_{s\in\Gamma_A}\beta(s)\#\{y:y^{\mathsf T}Z_s=0\}
 = \sum_{y\in\mathbf F_{p_0}^\ell}\frac{1}{p_0^\ell}
    \sum_{v\in\mathbf F_{p_0}^\ell}\ \sum_{s\in\Gamma_A}
        \psi(y^{\mathsf T}Z_s v)\beta(s).
\]
The total contribution from \(y=0\) or \(v=0\) is at most \(2Q\).
For \(y,v\ne0\), use the orthonormal basis of point masses and the
orthonormal basis of additive characters. Tensor the point mass at
\(\pi(s)A_0v\) and, separately, the normalized character at
\(\pi(s)^{-\mathsf T}A_1y\) with \(X_s\) and sum each over \(s\).
The squared norm of each sum is \(|\Gamma_A|\) times the sum of \(\beta\)
over the corresponding stabilizer. Expanding the inner product of the two
sums gives the innermost character sum above times \(|\Gamma_A|p_0^{-r/2}\).
Cauchy--Schwarz therefore bounds the real part of that character sum
by \(p_0^{r/2}b^{-6k}Q\). Together with \(p_0^{\ell+r/2}\le b^{6k}\)
this gives
\[
  \sum_{\det Z_s=0}\beta(s)
  \le p_0^{-1}\bigl(2+p_0^{\ell+r/2}b^{-6k}\bigr)Q
  \le 3b^{-10}Q .
\]
Allowing the conjugate place supplies the stabilizer condition for the
linear forms under the contragredient of \(\pi\).
\leanref{card_mul_sum_det_eq_zero_le}
The frame has length at most \(b\). Summing over the minors bounds the
contribution from the complement of the open reductions by \(b^{-4}Q\).
\leanref{sum_not_isOpenElement_le}

The open reductions avoid the scalar cosets, and \(r_\psi(s)=0\) there.
Apply \eqref{eq:span-cutoff} and \eqref{eq:span-open-input} with
\(D_{\mathrm{id}}\le b^{-6k|A|}D_{\mathrm{main}}\) and
\(b^{-15}\le\tfrac12 b^{-4}\). The square root of the contribution of the
open reductions is then at most
\((1+b^{-4})(b^{-6k|A|}D_{\mathrm{main}})^{1/2}\).
Combining the two contributions gives
\[
 Q \leq (1+4b^{-4})b^{-6k|A|}D_{\mathrm{main}}+b^{-4}Q,
 \qquad
 Q\leq (1+9b^{-4})b^{-6k|A|}D_{\mathrm{main}}.
\]
The second inequality follows from \((1-u)(1+9u)\ge1+4u\) at \(u=b^{-4}\).
It holds also when \(A\) consists of all the chosen places above \(p_0\).
Comparing the maximum to the value at \(A=\varnothing\) proves
\begin{equation}\label{eq:span-second-moment}
 \sum_{s\in\bar G}\|P_s\|_{\mathrm{HS}}^2
    \leq (1+9b^{-4})D_{\mathrm{main}}.
\end{equation}
\leanref{sum_secondMoment_le}

\subsection{From the second moment to a rank bound}

Put \(\epsilon=(8b^2)^{-1}\). Combining \eqref{eq:span-dimension-mass}
with \eqref{eq:span-second-moment} cancels the factor
\(\deg_B\,\mathrm{mass}_r\). Since \(9b^{-4}\le\epsilon\), we obtain
\begin{equation}\label{eq:span-moment-comparison}
 \frac{h}{|\bar G|}\sum_{s\in\bar G}\|P_s\|_{\mathrm{HS}}^2
    \leq (1+\epsilon)D_X^2.
\end{equation}
For every subspace \(U\subset S_0\), \eqref{eq:span-near-isometry}
also gives \(\|e_0P Q_U\|_{\mathrm{HS}}^2\geq\dim U-\epsilon D_X\).

Let \(F_1\subset V\) be a subspace and set \(\alpha=\dim F_1/h\).
Write \(Q_{F_1}\) for the orthogonal projection in \(\mathcal H\) onto \(F_1\).
The second moment gives
\begin{align*}
 \frac1{|\bar G|}\sum_{s\in\bar G}\|Q_{F_1}\rho(s)P\|_{\mathrm{HS}}^2
 &\leq
 \left(\dim F_1\cdot
        \frac1{|\bar G|}\sum_{s\in\bar G}\|P_s\|_{\mathrm{HS}}^2\right)^{1/2}\\
 &\leq \sqrt{\alpha(1+\epsilon)}D_X
 \leq \frac{1+\alpha+\epsilon}{2}D_X.
\end{align*}
The first line applies Cauchy--Schwarz to the Hilbert--Schmidt inner
product of \(Q_{F_1}\) with \(\sum_s X_s\). Indeed,
\(\|\sum_s X_s\|_{\mathrm{HS}}^2=|\bar G|\sum_s\beta(s)\).
Some \(s\) consequently attains a value no greater than this average.
Set \(B_s=e_0\rho(s)P\colon S_0\to V\) and \(U=B_s^{-1}(F_1)\).
On \(U\) the map \(B_s\) agrees with \(Q_{F_1}\rho(s)P\). Since \(\rho(s)\)
is unitary and commutes with \(e_0\), we have
\[
 \dim U-\epsilon D_X
 \ \leq\ \|e_0P Q_U\|_{\mathrm{HS}}^2
 \ \leq\ \|Q_{F_1}\rho(s)P\|_{\mathrm{HS}}^2.
\]
Therefore \(\dim U\le(1+\alpha+3\epsilon)D_X/2\).
\leanref{isSpread_of_second_moment}

The map to the quotient \(V/F_1\) has rank \(D_X-\dim U\).
While \(\operatorname{codim}F_1>16\epsilon h\), adjoining
\(\operatorname{im}B_s\) accounts for at least a fraction \(2D_X/(5h)\)
of the missing dimensions, since
\[
 \frac{1-\alpha-3\epsilon}{2}D_X
 \ge \frac25(1-\alpha)D_X
 \qquad (1-\alpha>16\epsilon).
\]
Write \(\lambda=2D_X/(5h)>0\). At \(F_1=0\) the gain is at most \(h\),
so \(\lambda\le1\). Continue choosing translates while the missing fraction
exceeds \(16\epsilon\). At each step the codimension decreases by a
factor at most \(1-\lambda\le e^{-\lambda}\). Within
\(\lceil\lambda^{-1}\log(1/(16\epsilon))\rceil\) steps it is at most
\(16\epsilon h=2h/b^2\). The set \(T\) of chosen translates satisfies
\[
  |T|D_X\le\tfrac52\bigl(1+\log(1/(16\epsilon))\bigr)h
               \le 5h\log(1+2b^2).
\]
\leanref{IsSpread.exists_few_le_finrank_iSup_range}

\subsection{Restriction to the ramified source}

Write \(N=\pi_N\overline{\pi_N}\), and for each \(w\mid\pi_N\) let \(z_w\)
be the finite adelic scalar action chosen in the formalization.
The space \(S\) consists of the elements killed by \(\sum_{j=0}^{N-2}z_w^j\)
for every \(w\). Since \(z_w^{N-1}=1\), the rank of this averaging operator
equals the dimension of the fixed space of \(z_w\).
The level generated by adjoining \(z_w\) has index at least \(N-1\) over
the original level. The class automorphism groups at both levels are
trivial, so the fixed space has dimension at most \(D_X/(N-1)\).
As there are \(f\) places \(w\mid\pi_N\), this gives
\[
 \operatorname{codim}_{S_0} S\le \frac{f}{N-1}D_X\le b^{-8}D_X.
\]
\leanref{finrank_ker_one_sub_scalingGenerator_le}

Restricting each map to \(S\) loses at most the dimension of the image
of a complement of \(S\). Since all the normalizing scalars are nonzero,
\(\operatorname{im}(B_s|_S)=\mathfrak p_sS\). The total loss is at most
\[
   |T|\operatorname{codim}_{S_0}S
   \le 5b^{-8}\log(1+2b^2)h\le 2h/b^2.
\]
Adding the earlier codimension bound \(2h/b^2\) gives a codimension
at most \(4h/b^2\). This and \(\dim S\le D_X\) prove the two inequalities
in Theorem~B. \leanref{enough_forms}
